\section{Eisenstein 级数}\srcpage{45}
\providecommand{\HH}{\mathbb{H}}% 本节新增宏，已在 reports/s10.md 登记

Langlands 理解自守形式空间的途径，基于一个归纳性的想法：先理解连续谱，\srcpage{46}而连续谱可以用 Levi 子群上的自守形式、借助 Eisenstein 级数来描述。回忆最经典的 Eisenstein 级数是如下求和
\[
E(\tau,s)=\frac{1}{2}\sum_{\substack{(m,n)\in\Z^{2}\\ \text{coprime}}}\frac{y^{s}}{|m+n\tau|^{2s}}
\]
其中 $\tau=x+iy$ 位于上半平面 $\HH=\{\tau\in\C\mid y=\operatorname{Im}(\tau)>0\}$ 内；该上半平面是 $\mathrm{SL}_2(\R)$ 经分式线性变换作用下的对称空间。而 $s\in\C$ 先验地要求 $\operatorname{Re}(s)>1$，以使该和收敛，但结果表明它可以亚纯延拓。由定义容易看出，$E(\tau,s)$ 下降为 $\HH/\mathrm{SL}_2(\Z)$ 上的函数。

Eisenstein 级数还有另一种写法。首先，
\[
E(\tau,s)=\sum_{\gamma\in\mathrm{SL}_2(\Z)/B(\Z)}\operatorname{Im}(\gamma(\tau))^{s},
\]
其中 $B(\Z)\subset\mathrm{SL}_2(\Z)$ 表示上三角矩阵构成的 Borel 子群（这里它就是 $\{\pm1\}\times\Z$）。事实上，稍作计算可知，对 $\tau=x+iy$ 有
\[
\operatorname{Im}(\gamma(\tau))=\operatorname{Im}\frac{a\tau+b}{c\tau+d}=\frac{y}{|c\tau+d|^{2}}.
\]

再作改写，我们可以考察由映射 $T\leftarrow B\to\mathrm{SL}_2$ 诱导的对应
\[
\begin{tikzcd}[column sep=large]
& B(\Z)\backslash B(\R)/\{\pm1\} \arrow[dl] \arrow[dr] & \\
T(\Z)\backslash T(\R)/\{\pm1\} & & \mathrm{SL}_2(\Z)\backslash\mathrm{SL}_2(\R)/\mathrm{SO}_2(\R)
\end{tikzcd}
\]
这里，右侧总取对极大紧子群的商。注意 $B(\R)/\{\pm1\}\to\mathrm{SL}_2(\R)/\mathrm{SO}_2(\R)$ 是同构。在此等同之下，投影 $B(\R)/\{\pm1\}\to T(\R)/\{\pm1\}=\R_{>0}$ 对应于投影
\[
\HH\to\R_{>0}\colon \tau=x+iy\mapsto\operatorname{Im}(\tau)=y.
\]
在这一解释下，我们从「$T=\Gm$ 上的自守形式」
\[
T(\Z)\backslash T(\R)/\{\pm1\}\cong\R_{>0}\to\C\colon y\mapsto|y|^{s}
\]
出发，把它拉回到 Borel 对应的空间 $B(\Z)\backslash B(\R)/\{\pm1\}$，再沿
\[
B(\Z)\backslash B(\R)/\{\pm1\}\to\mathrm{SL}_2(\Z)\backslash\mathrm{SL}_2(\R)/\mathrm{SO}_2(\R)
\]
的纤维求和，就得到 Eisenstein 级数 $E(\tau,s)$。

\srcpage{47}在这一解释下，显然可以看出如何把它推广得面目全非。用丛的模空间来说，我们可以考察对应
\[
\begin{tikzcd}[column sep=large]
& \Bun_B \arrow[dl,"p"'] \arrow[dr,"q"] & \\
\Bun_T & & \Bun_G,
\end{tikzcd}
\]
其中 $G$ 仍如前面，是带 Borel 子群 $B$ 与极大环面 $T$ 的分裂约化群。这诱导出函子
\[
\Eis_{!}=q_{!}p^{*}\colon D(\Bun_T)\to D(\Bun_G).
\]
实际上也可以把这个对应反向使用，从而得到常数项函子
\[
\CT_{!}=p_{!}q^{*}\colon D(\Bun_G)\to D(\Bun_T);
\]
在自守形式理论中，这对应于「沿尖点取 Fourier 展开中的常数项」。

更一般地，对 Levi 商为 $M$ 的抛物子群 $P$，我们得到对应
\[
\begin{tikzcd}[column sep=large]
& \Bun_P \arrow[dl,"p"'] \arrow[dr,"q"] & \\
\Bun_M & & \Bun_G.
\end{tikzcd}
\]
这给出类似的 Eisenstein 函子
\[
\Eis^{P}_{!}=q_{!}p^{*}\colon D(\Bun_M)\to D(\Bun_G)
\]
与常数项函子
\[
\CT^{P}_{!}=p_{!}q^{*}\colon D(\Bun_G)\to D(\Bun_M).
\]
显然，若把 $!$ 与 $*$ 分别换成 $*$ 与 $!$，可以得到这些函子彼此伴随的版本。有些出人意料的是，伴随性对上述这些函子本身也成立，只需把 $P$ 换成一个对立的抛物子群 $P^{-}$。按惯例这被称为「第二伴随」，因为它是 $p$ 进群表示论中 Bernstein 第二伴随的类比。\footnote{这一类比在 Hamann、Hansen 与 Scholze（2024）的设定中被处理得非常精确。}

\begin{theorem}[第二伴随；Drinfeld 与 Gaitsgory，2016]
函子 $\CT^{P^{-}}_{!}$ 是 $\Eis^{P}_{!}$ 的右伴随。该伴随关系在核的 $(\infty,2)$-范畴层面已经成立。
\end{theorem}

很容易给出该伴随的单位，即 $D(\Bun_M)$ 的自函子之间的一个映射
\[
\id\to\CT^{P^{-}}_{!}\circ\Eis^{P}_{!}.
\]
事实上，它由核的一个映射
\[
\Delta_{\Bun_M,!}\mathbf{1}\to f_{!}\mathbf{1}
\]
诱导。

\srcpage{48}这里 $f\colon\Bun_P\times_{\Bun_G}\Bun_{P^{-}}\to\Bun_M\times\Bun_M$ 是到两个 Levi 因子的投影。事实上，有开浸入
\[
*/M\hookrightarrow */P\times_{*/G}*/P^{-},
\]
它来自 $P\backslash G/P^{-}$ 中的开 Bruhat 胞腔 $*/M$；对 $X$ 取映射叠之后，这诱导开嵌入
\[
j\colon\Bun_M\to\Bun_P\times_{\Bun_G}\Bun_{P^{-}}.
\]
于是把 $f_{!}$ 作用于 $j_{!}\mathbf{1}\to\mathbf{1}$，即得所求的映射。

Drinfeld 与 Gaitsgory 在核的层面上给出该伴随余单位的一个几何构造；这用到一些与双曲局部化有关的想法，以及一族位于 $\A^{1}/\Gm$ 上的群 $\tilde G$，其一般纤维是 $G$，特殊纤维是 $P\times_M P^{-}$。\footnote{这一族 $\tilde G$ 也可以作为 Vinberg 幺半群的某一部分得到。}更确切地说，我们有一族
\[
\tilde G\to G\times G\times\A^{1}/\Gm,
\]
它诱导
\[
\Bun_{\tilde G}\to\Bun_G\times\Bun_G\times\A^{1}/\Gm,
\]
在对角嵌入与给出复合 $\Eis^{P}_{!}\circ\CT^{P^{-}}_{!}$ 的循环 $\Bun_P\times_{\Bun_M}\Bun_{P^{-}}$ 之间插值。借助一个余特殊化映射，由此可以造出所求的映射。

Eisenstein 级数的一个重要方面是：它们与 Hecke 算子大体上交换，从而把 Hecke 特征层送到 Hecke 特征层。严格说来，这并不完全成立。尽管如此，Campbell、Chen、Færgeman、Gaitsgory、Lin、Raskin 与 Rozenblyum（2024）在 D-模情形证明了下图差一个扭交换：
\[
\begin{tikzcd}[column sep=large,row sep=large]
D(\Bun_M) \arrow[r,"\mathbb{L}_M"] \arrow[d,"\Eis^{P}_{!}"'] & \IndCoh_{\Nilp}(\LocSys_{\Md}) \arrow[d,"\Eis^{P}_{\mathrm{spec}}"] \\
D(\Bun_G) \arrow[r,"\mathbb{L}_G"] & \IndCoh_{\Nilp}(\LocSys_{\Gd}).
\end{tikzcd}
\]
这里，谱 Eisenstein 函子指谱侧的一个类似构造。其证明用到：用 Whittaker 范畴 $\Whit(G)$ 对 Whittaker 层上的谱作用所作的描述，以及几何 Casselman--Shalika 公式与常数项的相容性。

\begin{remark}
Eisenstein 级数归根结底正是必须把 $D_{\qc}(\LocSys_{\Gd})$ 换成 $\IndCoh_{\Nilp}(\LocSys_{\Gd})$ 的原因。也就是说，函子 $\Eis^{P}_{!}$ 保持紧对象；但若把它的谱侧对应物看作 $D_{\qc}$ 上的函子，则后者不保持紧对象。真前推保持凝聚层的复形，却不保持完美复形。\srcpage{49}尽管如此，此时引入的奇异支撑方向必须沿给定抛物子群 $P$ 的单幂根而是幂零的。
\end{remark}

在上一个图中竖直地取右伴随，得到另一个图
\[
\begin{tikzcd}[column sep=large,row sep=large]
D(\Bun_M) \arrow[r,"\mathbb{L}_M"] & \IndCoh_{\Nilp}(\LocSys_{\Md}) \\
D(\Bun_G) \arrow[r,"\mathbb{L}_G"] \arrow[u,"\CT^{P^{-}}_{!}"] & \IndCoh_{\Nilp}(\LocSys_{\Gd}) \arrow[u,"\CT^{P^{-}}_{\mathrm{spec}}"']
\end{tikzcd}
\]
它先验地只到差一个自然归一化为止交换。Campbell、Chen、Færgeman、Gaitsgory、Lin、Raskin 与 Rozenblyum（2024）的一个主要定理是：这个图确实交换。其证明用到 Kac--Moody 局部化。
